
\documentclass[openany, amssymb, psamsfonts]{amsart}
\usepackage{mathrsfs,comment}
\usepackage[usenames,dvipsnames]{color}
\usepackage[normalem]{ulem}
\usepackage{url}
\usepackage[all,arc,2cell]{xy}
\UseAllTwocells
\usepackage{enumerate}
\usepackage{hyperref}
\hypersetup{%
  bookmarksnumbered=true,%
  bookmarks=true,%
  colorlinks=true,%
  linkcolor=blue,%
  citecolor=blue,%
  filecolor=blue,%
  menucolor=blue,%
  pagecolor=blue,%
  urlcolor=blue,%
  pdfnewwindow=true,%
  pdfstartview=FitBH}

\let\fullref\autoref

\def\makeautorefname#1#2{\expandafter\def\csname#1autorefname\endcsname{#2}}

\def\equationautorefname~#1\null{(#1)\null}
\makeautorefname{footnote}{footnote}%
\makeautorefname{figure}{Figure}%
\makeautorefname{table}{Table}%
\makeautorefname{section}{Section}%
\makeautorefname{subsection}{Section}%
\makeautorefname{theorem}{Theorem}%
\makeautorefname{thm}{Theorem}%
\makeautorefname{cor}{Corollary}%
\makeautorefname{lem}{Lemma}%
\makeautorefname{prop}{Proposition}%
\makeautorefname{defn}{Definition}%
\makeautorefname{rem}{Remark}%
\makeautorefname{obs}{Observation}%

\newtheorem{thm}{Theorem}[section]
\newtheorem{cor}{Corollary}[section]
\newtheorem{prop}{Proposition}[section]
\newtheorem{lem}{Lemma}[section]

\theoremstyle{definition}
\newtheorem{defn}{Definition}[section]
\newtheorem{exmp}{Example}[section]
\newtheorem{rem}{Remark}[section]
\newtheorem{obs}{Observation}[section]

\makeatletter
\let\c@cor=\c@thm
\let\c@prop=\c@thm
\let\c@lem=\c@thm
\let\c@defn=\c@thm
\let\c@exmp=\c@thm
\let\c@rem=\c@thm
\let\c@obs=\c@thm
\let\c@equation\c@thm
\numberwithin{equation}{section}
\makeatother

\usepackage{booktabs}
\usepackage{multirow}

\bibliographystyle{plain}

\title{Neural Theorem Proving on miniF2F: \\
Scaling, Tree Search, and the Limits of Fine-Tuning}

\author{Aslan Feng}

\date{Columbia Business School REU, Summer 2026 \\
Mentor: Professor Namkoong}

\begin{document}

\begin{abstract}
We investigate how far a single off-the-shelf theorem-proving model can be pushed on
the miniF2F benchmark under realistic compute constraints: one A40 GPU, 4-bit
quantization, and a fixed per-problem time budget.
Whole-proof generation with DeepSeek-Prover-V1.5-RL reaches 39.8\% pass@1 at
64 samples per problem, then plateaus completely: scaling to 256 samples yields
zero additional proofs.
Adding Monte Carlo Tree Search (BFS, depth at most 6, verified via a single
\texttt{lake build} call) improves the score to 48.8\%, with tree search uniquely
closing 25 problems that exhaustive whole-proof sampling never finds.
Two attempts at LoRA fine-tuning both degraded performance relative to the
unmodified base model.
We trace this failure to a distributional mismatch: 98\% of the available
fine-tuning proofs are single-line tactic closers, while 63\% of the problems
successfully proved on miniF2F require multi-step reasoning.
The remaining gap to the published 60.2\% result reflects a combination of
compute scale (400 vs.\ 3200 samples per problem) and a hard tail of
122 competition problems that neither strategy solves.
\end{abstract}

\maketitle

\tableofcontents

%-----------------------------------------------------------------------
\section{Introduction}
%-----------------------------------------------------------------------

Automated theorem proving has long been a benchmark for machine reasoning.
The introduction of large language models trained specifically on formal
mathematics---in particular, on Lean 4 and its Mathlib library---has led
to dramatic improvements on standard benchmarks.
DeepSeek-Prover-V1.5-RL \cite{deepseek} achieves 60.2\% on the miniF2F
benchmark \cite{minif2f}, which consists of 244 competition mathematics problems
formalized in Lean 4, using reinforcement learning from proof feedback and
Monte Carlo Tree Search with thousands of samples per problem.

This project asks a more constrained question: how much of that performance
can be recovered with a single GPU node, 4-bit quantized inference, and
roughly 400 samples per problem?
We find the answer to be 48.8\%---well above the 24.6\% achieved by simpler
baselines, but still 11 percentage points short of the paper's result.
The analysis of where that gap comes from, and what happened when we tried
to close it, turns out to be as informative as the headline number.

Three main findings drive the paper.
First, whole-proof sampling saturates very early: going from 64 to 256 samples
per problem yields no new proofs, indicating that the 60\% barrier requires
a fundamentally different strategy rather than more of the same computation.
Second, tree search does help, but asymmetrically: it solves 25 problems that
exhaustive sampling never finds (mostly harder multi-step competition problems)
while losing 3 problems that sampling gets lucky on, for a net gain of
22 additional theorems.
Third, fine-tuning the model on available Lean 4 proof corpora
consistently makes things worse, because nearly all available training
proofs are single-line tactic applications while the miniF2F problems that
need help are precisely the ones requiring multi-step structured reasoning.

%-----------------------------------------------------------------------
\section{Background}
%-----------------------------------------------------------------------

\subsection{The miniF2F Benchmark}

miniF2F \cite{minif2f} is a curated collection of 488 competition mathematics
problems (244 validation, 244 test) drawn from AMC, AIME, and IMO competitions,
formalized in Lean 4.
Each problem provides a complete Lean 4 theorem declaration ending in
\texttt{:= sorry}, and the task is to replace \texttt{sorry} with a valid proof.

The problems span three main categories: algebra (88 problems, 36\% of the
test set), number theory (68 problems, 28\%), and competition-style problems
from AMC/AIME/IMO (80 problems, 33\%), with the remainder being miscellaneous
mathematical statements.
Difficulty varies widely, from one-line arithmetic verifications to
multi-step inequalities requiring specialized Mathlib lemmas.

\subsection{DeepSeek-Prover-V1.5-RL}

DeepSeek-Prover-V1.5-RL \cite{deepseek} is a 7-billion-parameter causal language
model derived from DeepSeek-Coder-V1.5-Instruct and further trained via reinforcement
learning on Lean 4 proof feedback.
Its inference protocol is \emph{whole-proof generation}: the model is given a
Lean 4 file stub ending with \texttt{:= by}, and generates the complete proof body
in a single forward pass.
The published result of 60.2\% on miniF2F-test uses roughly 3200 independent
samples per problem and a tree-search variant that applies at each step.

We use the publicly released weights under 4-bit NF4 quantization via
BitsAndBytes \cite{bitsandbytes}, reducing memory from approximately 14 GB to
4 GB and enabling the full model to run on a single A40 GPU alongside the
Lean verification infrastructure.

\subsection{Lean 4 Verification}

Each candidate proof is verified by compiling it against the Mathlib library
using \texttt{lake build} \cite{mathlib}.
A proof is accepted if and only if the build exits with code 0 and no
\texttt{uses 'sorry'} warning appears in the output.
For tree-search verification, all candidates at each depth are compiled in a
single batch file, with errors attributed to individual theorems by line range.
For whole-proof verification, up to 4 candidate files are compiled in parallel
to avoid exhausting GPU-node memory (each Lean elaboration process loads the
full Mathlib olive olean cache, approximately 8--12 GB per process).

%-----------------------------------------------------------------------
\section{Methods}
%-----------------------------------------------------------------------

\subsection{Whole-Proof Sampling}

The baseline strategy generates $k$ independent proof attempts per problem via
temperature sampling ($T = 1.0$, $p = 0.95$), compiles them against Lean 4 via
\texttt{lake build}, and reports success if any candidate passes.
We evaluate at $k \in \{32, 64, 256\}$, with a 300-second timeout per
verification call and a per-proof token budget of 256 new tokens.

All generation runs use batch size 16, which we found to be the largest
stable batch size under 4-bit quantization.
Batch size 8 produces deadlocks in \texttt{torch.multinomial} on CUDA 12.2 due
to a known interaction with BitsAndBytes dequantization; this required careful
handling throughout the evaluation.

\subsection{Monte Carlo Tree Search}

The tree-search variant explores proof tactics step-by-step using BFS to
depth at most 6, with width 8 (eight tactic candidates generated per
node expansion).
At each depth, all candidate tactic sequences are compiled in a single
\texttt{lake build} call on a shared \texttt{ProofGoals.lean} file.
Errors are attributed to individual theorems by tracking each theorem's
line range in the file.

One subtle correctness issue required a fix: Lean sometimes reports
``unsolved goals'' errors at the blank separator line after the last
tactic rather than on the tactic itself (at block-close elaboration time).
The naive implementation did not include this separator line in the range
assigned to a theorem, causing roughly 51 false positives in which
incomplete proofs appeared to succeed.
The fix extends each theorem's range by one line to capture separator-line
errors.

Additionally, a tactic filter drops any candidate that matches
\verb|^\s*theorem\s+\S|, which catches cases where the model regenerates
the theorem header as a tactic---an error that caused Lean to silently
close the current \texttt{by} block early, producing roughly 17 additional
false positives in the uncorrected implementation.

Each problem receives a 120-second tree-search budget followed by a
whole-proof fallback phase using 400 samples total.
If tree search proves the goal before the budget expires, the fallback
is skipped.

\subsection{LoRA Fine-Tuning}

We conducted two rounds of LoRA fine-tuning on the DeepSeek-Prover base model,
both using whole-proof training pairs (Lean 4 file stub, complete proof body)
to match the inference format.

\textbf{Version 1} trained on 67 examples: the validated miniF2F proofs found
by the base model on the validation split.
This run used 4-bit NF4 quantization with float16 compute and LoRA rank 16.

\textbf{Version 2} addressed the data scarcity by adding 10,434 proofs from the
Lean-Workbook dataset \cite{leanworkbook} (filtered to examples with
\texttt{state\_after = "no goals"}), giving 10,561 total training examples.
This run used full BF16 precision and LoRA rank 64 on an A100-class GPU.

%-----------------------------------------------------------------------
\section{Results}
%-----------------------------------------------------------------------

\autoref{tab:main_results} summarizes all evaluated configurations on the
miniF2F test split (244 problems).
All results with an asterisk are confirmed by single-candidate re-verification
using \texttt{lake env lean} on a GPU node; the remainder are verified by
the batch \texttt{lake build} approach described in \autoref{sec:analysis_mcts}.

\begin{table}[h]
\centering
\caption{miniF2F-test results (244 problems). ``Verified'' means the reported
score was confirmed by independent re-verification to exclude batch false positives.}
\label{tab:main_results}
\begin{tabular}{@{}llccc@{}}
\toprule
Model & Method & Proved & Pass@1 & Verified \\
\midrule
ByT5-small pretrained & tactic + fallback, $k=32$ & 59 & 24.2\% & $\star$ \\
ByT5-small analysis FT & tactic + fallback, $k=32$ & 60 & 24.6\% & $\star$ \\
ByT5-small Mathlib FT (ep.\,5) & tactic + fallback, $k=32$ & 70 & 28.7\% & $\star$ \\
DeepSeek base & whole-proof, $k=32$ & 60 & 24.6\% & $\star$ \\
DeepSeek + LoRA V1 & whole-proof, $k=32$ & 45 & 18.4\% & $\star$ \\
DeepSeek + LoRA V2 & whole-proof, $k=32$ & 41 & 16.8\% & $\star$ \\
DeepSeek base & whole-proof, $k=64$ & 97 & 39.8\% & $\star$ \\
DeepSeek base & whole-proof, $k=256$ & 97 & 39.8\% & $\star$ \\
\textbf{DeepSeek + MCTS} & \textbf{tree search + $k=400$} & \textbf{119} & \textbf{48.8\%} & $\star$ \\
\midrule
ReProver \cite{reprover} & best-first search & -- & $\approx$26.5\% & -- \\
DeepSeek-Prover-V1.5-RL \cite{deepseek} & MCTS, $k \approx 3200$ & 147 & 60.2\% & -- \\
\bottomrule
\end{tabular}
\end{table}

The MCTS result of 48.8\% is our primary contribution.
It is fully verified via the corrected batch lake build verifier, and
represents a 9.0 percentage-point improvement over the whole-proof
sampling ceiling of 39.8\%.
The ByT5 results are included for completeness; they are discussed briefly
in \autoref{sec:byt5}.

%-----------------------------------------------------------------------
\section{Analysis}
\label{sec:analysis}
%-----------------------------------------------------------------------

\subsection{Sample Scaling Saturates Early}
\label{sec:scaling}

\autoref{tab:scaling} shows the marginal gain from adding more whole-proof
samples.

\begin{table}[h]
\centering
\caption{Marginal gains from increasing the sample budget.}
\label{tab:scaling}
\begin{tabular}{@{}lccr@{}}
\toprule
Run & Proved & Pass@1 & New proofs \\
\midrule
$k = 32$ & 60 & 24.6\% & -- \\
$k = 64$ & 97 & 39.8\% & $+$37 \\
$k = 256$ & 97 & 39.8\% & $\phantom{+}$0 \\
MCTS ($k \approx 400$) & 119 & 48.8\% & $+$22 vs.\ $k=256$ \\
\midrule
Published ($k \approx 3200$) & 147 & 60.2\% & -- \\
\bottomrule
\end{tabular}
\end{table}

Going from 32 to 64 samples gains 37 theorems---a large jump that reflects
the model finding correct proofs for many problems on the second or third
try.
Going from 64 to 256 samples (four times more computation) gains nothing.
This is not a coincidence of rounding: of the 4 problems that the
verified top-256 result claimed to prove beyond top-64,
all 4 turned out to be false positives from the original batch verifier.

The implication is stark: additional sampling, in isolation, cannot close
the gap to 60.2\%.
Any path from 39.8\% to 60\% requires a qualitatively different strategy.
MCTS provides part of that---contributing 22 genuinely new proofs---but
the remaining 28 problems (the gap between our 48.8\% and the paper's 60.2\%)
are unlikely to yield to more of the same.

The shape of the scaling curve (a jump, then a hard plateau) suggests that
the model's sampling distribution covers a fixed ``easy'' region of
proof space, and the hard tail of 122 remaining problems lies outside
that region entirely regardless of how many samples are drawn.

\subsection{Tree Search vs.\ Whole-Proof Sampling}
\label{sec:analysis_mcts}

\autoref{tab:overlap} compares MCTS and $k=256$ sampling problem-by-problem.

\begin{table}[h]
\centering
\caption{Problem-level overlap between MCTS and whole-proof sampling ($k=256$).
Both methods were applied to all 244 test problems.}
\label{tab:overlap}
\begin{tabular}{@{}lcc@{}}
\toprule
Outcome & Problems & Fraction \\
\midrule
Both methods proved & 94 & 38.5\% \\
Only MCTS proved & 25 & 10.2\% \\
Only $k=256$ proved & 3 & 1.2\% \\
Neither proved & 122 & 50.0\% \\
\midrule
\textit{Total} & 244 & \\
\bottomrule
\end{tabular}
\end{table}

MCTS uniquely solves 25 problems that 256 whole-proof attempts never find.
Inspection of these problems reveals that most are multi-step: their
MCTS proofs average 8.2 tactic steps (median 5), compared to 3.8 steps
(median 2) for problems both methods solve.
Tree search succeeds here by constructing the proof incrementally:
it verifies each tactic against Lean before committing, pruning
invalid branches and building toward a complete proof step-by-step.
Whole-proof sampling must generate the entire correct proof in one shot,
which becomes exponentially unlikely as the number of required steps grows.

The 3 problems solved only by $k=256$ but not by MCTS are instructive in
the opposite direction.
All three have proofs that whole-proof sampling stumbled upon during
its 256 attempts, but whose first tactic did not survive MCTS's pruning
(either because it was rare in the model's distribution, or because
the particular formulation clashed with how the batch verifier
attributes errors).
This is a known limitation of the batch-verifier approach: at depth 1
of BFS, the 8 generated tactics must include the right starting tactic,
and if the model assigns that tactic low probability, it will be missed.
A broader initial width would recover these cases at the cost of more
verification time.

Overall, the two strategies are complementary rather than redundant:
MCTS earns its keep on problems where structured step-by-step construction
matters, while whole-proof sampling is occasionally more effective for
problems with one-shot ``lucky'' proofs.

\subsection{Category-Level Breakdown}
\label{sec:categories}

\autoref{tab:categories} breaks down performance by problem type.

\begin{table}[h]
\centering
\caption{Per-category results on miniF2F-test.}
\label{tab:categories}
\begin{tabular}{@{}lcccc@{}}
\toprule
Category & Total & $k=256$ & MCTS & MCTS gain \\
\midrule
Algebra & 88 & 52\,(59\%) & 58\,(66\%) & $+$6 \\
Number theory & 68 & 29\,(43\%) & 32\,(47\%) & $+$3 \\
Competition (AMC/AIME/IMO) & 80 & 13\,(16\%) & 26\,(33\%) & $+$13 \\
\midrule
All & 244 & 97\,(39.8\%) & 119\,(48.8\%) & $+$22 \\
\bottomrule
\end{tabular}
\end{table}

Algebra problems are the most amenable to both strategies (59\%/66\%),
likely because many can be resolved with a small set of universally
applicable tactics: \texttt{linarith}, \texttt{ring}, \texttt{norm\_num},
and their combinations.
Number theory is harder (43\%/47\%) despite often having short formal statements,
because the proofs tend to require modular arithmetic reasoning that
the model must produce precisely.

Competition problems are where the two strategies diverge most.
Whole-proof sampling solves only 16\% (13/80)---these are the
``lucky'' problems where the model happens to generate the right proof
in 256 tries.
MCTS doubles this to 33\% (26/80) by finding multi-step solutions
that sampling would never produce.
Nevertheless, 67\% of competition problems (54/80) remain unsolved by
either method.
These are the genuinely hard problems: they require domain-specific
Mathlib lemmas, intricate \texttt{have} chains, or calculation blocks
that neither whole-proof generation nor 6-step tree search reaches.

The category gap also reveals where additional samples would be most
wasted: the algebra and number theory problems already have high success
rates, so more computation yields diminishing returns on those categories.
Any meaningful path to 60\% must come from the competition category,
which requires either a fundamentally better model or a longer-horizon
search strategy.

\subsection{Fine-Tuning Degrades Whole-Proof Generation}
\label{sec:lora}

Both LoRA fine-tuning attempts degraded performance relative to the
unmodified base model, despite using progressively better configurations:

\begin{table}[h]
\centering
\caption{LoRA fine-tuning results compared to base model.}
\label{tab:lora}
\begin{tabular}{@{}llccc@{}}
\toprule
Model & Training data & Proved & Pass@1 & $\Delta$ vs.\ base \\
\midrule
DeepSeek base & -- & 60 & 24.6\% & -- \\
DeepSeek + LoRA V1 & 67 miniF2F proofs & 45 & 18.4\% & $-$6.2 pp \\
DeepSeek + LoRA V2 & 10,561 proofs & 41 & 16.8\% & $-$7.8 pp \\
\bottomrule
\end{tabular}
\end{table}

LoRA V1 used 67 training examples (validated miniF2F proofs) with 4-bit
quantization and LoRA rank 16; it scored 18.4\%, losing 6.2 pp.
LoRA V2 addressed data scarcity with 10,561 training examples and full BF16
precision, using LoRA rank 64; it scored 16.8\%, losing an additional 1.6 pp.
More data and better precision did not help---the degradation actually
increased slightly.

The root cause is a distributional mismatch between training and inference.
When we examine the Lean-Workbook corpus used for LoRA V2, we find:

\begin{itemize}
\item 98.0\% of training proofs are single-line (one tactic closes the goal)
\item 0.8\% require exactly two lines
\item 1.2\% require three or more lines
\end{itemize}

The base model's training distribution (via reinforcement learning on Lean 4)
is broadly calibrated to the full difficulty range of mathematical theorem proving.
After fine-tuning on a corpus that is 98\% single-step proofs, the model's
output distribution shifts strongly toward generating short terminating tactics.
This is precisely the wrong inductive bias for miniF2F: among the 119 problems
successfully proved by MCTS, 63\% (75/119) require proofs with more than one step.
The fine-tuned model generates correct single-step terminations with high probability
but has learned to truncate multi-step proofs prematurely.

This finding is consistent with the LoRA V1 training report, which noted that
all 41 proofs found by LoRA V2 used short closers (\texttt{omega}, \texttt{linarith},
\texttt{ring}, \texttt{norm\_num}, \texttt{rfl})---identical to what the unmodified
base model finds greedily, with none of the more structured multi-step proofs
that distinguish the harder miniF2F problems.

\begin{obs}
Fine-tuning a whole-proof generation model on tactic-level or single-step
training data is counterproductive for benchmarks that require multi-step
reasoning. The appropriate training format is (file stub, complete proof body)
pairs drawn from the same difficulty distribution as the target benchmark.
\end{obs}

\subsection{Why the Gap to 60\% Remains}
\label{sec:gap}

The published 60.2\% result differs from ours in two dimensions beyond model
architecture and training:

\textbf{Sample count.} The paper uses approximately 3200 samples per problem;
we use 400. Our scaling analysis (\autoref{sec:scaling}) shows that doubling
or quadrupling the sample count within the whole-proof regime produces no gain
once the plateau is reached. However, the published system applies tree search
\emph{at every node} and generates new samples at each depth, effectively
distributing the 3200-sample budget across the search tree rather than
spending it all on independent full-proof attempts. This is a qualitatively
different use of compute that we do not replicate.

\textbf{The hard tail.} Of the 244 test problems, 122 (50\%) are solved by
neither method we evaluate. These problems are concentrated in the
competition category (54/80) and represent problems where both the proof
requires uncommon Mathlib lemmas and the number of required steps exceeds
our tree search depth of 6. The published system solves approximately
$147 - 119 = 28$ of these---roughly 23\%---suggesting that a deeper
or wider search, possibly combined with a better-calibrated model, could
recover a significant fraction.

%-----------------------------------------------------------------------
\section{The ByT5 Tactic Model}
\label{sec:byt5}
%-----------------------------------------------------------------------

As a secondary experiment, we evaluated ByT5-small \cite{byt5}
(300M parameters, byte-level tokenizer), fine-tuned on Mathlib tactic data
from LeanDojo \cite{leandojo}.
ByT5 generates one tactic at a time from a proof state, rather than a complete proof.
On miniF2F, it is augmented with 14 universal fallback tactics
(\texttt{norm\_num}, \texttt{omega}, \texttt{ring}, \texttt{linarith}, \texttt{decide},
\texttt{aesop}, and combinations).

After fine-tuning on the full Mathlib tactic dataset for 5 epochs
(${\approx}250\text{K}$ training examples), ByT5 reaches 70/244 = 28.7\%, surpassing
the published ReProver result of ${\approx}26.5\%$.
However, this performance is almost entirely explained by the 14 fallback tactics,
which mechanically close most of the straightforward miniF2F problems.
The fine-tuned model's tactic predictions contribute marginally beyond the pretrained
baseline: both pretrained and analysis-fine-tuned ByT5 achieve around 24\%,
with the gains from more training data being modest.

This contrasts with DeepSeek, where the 4-bit base model dramatically outperforms
ByT5 under the same sample budget once the whole-proof format is used correctly.
The two architectures probe fundamentally different capabilities: ByT5 is a precise
tactic predictor that requires an interactive REPL to apply tactics step by step,
while DeepSeek's causal LM format enables it to generate entire multi-step proofs
in one shot.

%-----------------------------------------------------------------------
\section{Implementation Challenges}
%-----------------------------------------------------------------------

Several engineering issues complicated the evaluation in ways that are worth
documenting for reproducibility.

\textbf{Verification false positives.}
The initial batch-verifier implementation produced a 47.6\% false positive rate
(68 of 143 claimed MCTS proofs were incorrect).
Two bugs were responsible: a line-range attribution error that missed
``unsolved goals'' errors reported on separator lines (51 false positives),
and a tactic filter that failed to drop cases where the model
regenerated the theorem header (17 false positives).
Both are fixed in the final implementation.

\textbf{Thread limits on CPU nodes.}
Lean spawns threads for parallel elaboration.
CPU cluster nodes impose strict thread limits that caused SIGABRT
when reverification was run on them.
All reverification was moved to GPU nodes, and
\texttt{LEAN\_NUM\_THREADS=1} was added to the subprocess environment
as a belt-and-suspenders fix.

\textbf{Memory constraints.}
Running 8 parallel \texttt{lake env lean} processes, each loading the
full Mathlib olean cache (${\approx}8$--12 GB per process), exhausted
100 GB nodes.
The whole-proof verification pipeline was capped at 4 parallel workers.

\textbf{BPE decode artifacts.}
DeepSeek's tokenizer produces Unicode byte-pair symbols (\texttt{Ġ}, \texttt{Ċ})
that must be post-processed to spaces and newlines; without this, all
generated proofs are garbled and the model appears to achieve 0\%.

\textbf{CUDA deadlock at batch size 8.}
A known bug in BitsAndBytes dequantization on CUDA 12.2 causes
\texttt{torch.multinomial} to deadlock at batch size 8.
Only batch size 16 (larger, counterintuitively) was stable.

%-----------------------------------------------------------------------
\section{Conclusion}
%-----------------------------------------------------------------------

We evaluated DeepSeek-Prover-V1.5-RL on miniF2F under constrained compute,
reaching 48.8\% with MCTS tree search---9 percentage points above the
whole-proof sampling ceiling and 22 above simpler baselines.
The main lessons are:

\begin{enumerate}
\item \textbf{Whole-proof sampling plateaus quickly.} The jump from 32 to 64
samples is large (+37 problems); from 64 to 256 is zero. More sampling
alone cannot close the gap to 60\%.
\item \textbf{Tree search and whole-proof sampling are complementary.}
MCTS uniquely solves 25 harder multi-step problems; sampling uniquely solves
3 ``lucky'' one-shot problems. Together they reach 48.8\% vs.\ 39.8\% for
sampling alone.
\item \textbf{Fine-tuning on the wrong distribution actively hurts.}
When 98\% of training proofs are single-step closers but the target benchmark
requires multi-step reasoning, fine-tuning suppresses exactly the capability
that is needed. Training data must match the inference difficulty profile.
\item \textbf{Competition problems are the hard tail.} Even with MCTS, only
33\% of AMC/AIME/IMO problems are solved. This category is where the gap to
60\% lives, and closing it likely requires deeper search, broader model
sampling, or better-aligned fine-tuning data.
\end{enumerate}

The results suggest that the published 60.2\% is achievable with sufficient
compute (more samples per tree-search node) and the right search budget,
but is not reachable by scaling up the simple strategies evaluated here.
This points toward richer proof-search algorithms and better training corpora
as the more promising directions for future work.

%-----------------------------------------------------------------------
\section*{Acknowledgments}
%-----------------------------------------------------------------------

It is a pleasure to thank Professor Namkoong for guidance throughout this
project and for feedback that shaped the analysis reported here.
Computations were carried out on the CBS Research Grid at Columbia Business
School.

%-----------------------------------------------------------------------
\begin{thebibliography}{9}

\bibitem{deepseek}
DeepSeek-AI.
DeepSeek-Prover-V1.5: Harnessing Proof Assistant Feedback for Reinforcement
Learning and Monte-Carlo Tree Search.
\textit{arXiv preprint arXiv:2408.08152}, 2024.

\bibitem{minif2f}
Kunhao Zheng, Jesse Michael Han, and Stanislas Polu.
miniF2F: a cross-system benchmark for formal Olympiad-level mathematics.
\textit{International Conference on Learning Representations}, 2022.

\bibitem{mathlib}
The Mathlib Community.
The Lean 4 Mathematical Library.
\url{https://leanprover-community.github.io/mathlib4_docs/}, 2024.

\bibitem{reprover}
Kaiyu Yang, Aidan Swope, Alex Gu, Rahul Chalamala, Peiyang Song, Shixing Yu,
Saad Godil, Ryan Prenger, and Anima Anandkumar.
LeanDojo: Theorem Proving with Retrieval-Augmented Language Models.
\textit{Advances in Neural Information Processing Systems}, 2023.

\bibitem{byt5}
Linting Xue, Aditya Barua, Noah Constant, Rami Al-Rfou, Sharan Narang,
Mihir Kale, Adam Roberts, and Colin Raffel.
ByT5: Towards a Token-Free Future with Pre-trained Byte-to-Byte Models.
\textit{Transactions of the Association for Computational Linguistics}, 2022.

\bibitem{leandojo}
Kaiyu Yang, Aidan Swope, Alex Gu, Rahul Chalamala, Peiyang Song, Shixing Yu,
Saad Godil, Ryan Prenger, and Anima Anandkumar.
LeanDojo: Theorem Proving with Retrieval-Augmented Language Models.
\url{https://github.com/lean-dojo/LeanDojo}, 2023.

\bibitem{bitsandbytes}
Tim Dettmers, Artidoro Pagnoni, Ari Holtzman, and Luke Zettlemoyer.
QLoRA: Efficient Finetuning of Quantized LLMs.
\textit{Advances in Neural Information Processing Systems}, 2023.

\bibitem{leanworkbook}
Huaiyuan Ying, Zijian Wu, Yihan Geng, Jiayu Wang, Dahua Lin, and Bo Dai.
Lean Workbook: A Large-Scale Lean Problem Set Formalized from Natural Language
Math Problems.
\textit{arXiv preprint arXiv:2406.03847}, 2024.

\end{thebibliography}

\end{document}
